\documentclass[../candidatura.tex]{subfiles}
\graphicspath{{\subfix{../imagenes/}}}
\begin{document}

\section{Resumen}

Este proyecto de tesis es una investigación artística que propone la conceptualización, práctica y materialización de la popusintesíntesis, entendida como la síntesis popular de sintetizadores, y en particular sus aplicaciones en sintetizadores sonoros.

El neologismo popusintesíntesis es una respuesta y tributo a ``synthesynthesis'', palabra propuesta por el artista y empresario estadounidense Peter Blasser en su tesis de magíster, donde define su práctica de diseño y fabricación de sintetizadores como sintesíntesis \footcite{blasser_stores_2005}. Su aproximación crítica es una respuesta a la industria de la música de producción de conciertos y grabaciones en distintos formatos, las corporaciones que empaquetan circuitos como cajas negras y las convierten en sintetizadores ofuscados con parámetros. La sintesíntesis de Peter Blasser emerge como una respuesta donde el valor de la obra artística está en el diseño y fabricación de sintetizadores, celebrando sus soldaduras, sus aserrines, sus imperfecciones.

El prefijo popular de la popusintesíntesis es otro homenaje y guiño a la obra de Tatsuya Takahashi, diseñador de sintetizadores. En su primera era trabajando en la empresa japonesa KORG entre 2010 y 2017 como ingeniero jefe de sintetizadores análogos, publicó un post en Facebook donde celebraba que los Korg volca fueran el primer sintetizador de infancias aprendiendo música, y que eso resonaba con el nombre volca, en alusión a la palabra alemana Volk, y a otras marcas como Volkswagen, el auto del pueblo. Tatsuya propone a los Korg volca como sintetizadores de la gente, populares \footcite{synthtopia_takahashi_2017}. En su visión, lo popular no termina en la compra del instrumento: sus diseños incluyen guiños a quienes quieran abrirlos y modificarlos, como la publicación de los esquemáticos del Korg monotron y las serigrafías en sus placas de circuito impreso que invitan a soldar nuevas conexiones \footcite{korg_monotron_sinfecha}.

La componente popular entonces es un signo sobrecargado de una constelación de definiciones que se abordan en esta tesis, para situar esta práctica dentro de la cultura popular, los objetos cotidianos, la política de software y hardware de fuente abierta, la reparabilidad, la legibilidad, y la aspiración de convertir la sintesíntesis en una práctica masiva de goce sonoro.

En esta tesis se proponen distintas estrategias para la síntesis popular de sintetizadores sonoros, con un fuerte énfasis en herramientas electrónicas y computacionales. La elección de herramientas usadas y propuestas también responde a lo popular en lo valórico, cristalizándose en software y hardware de fuente abierta, legibilidad por sobre eficiencia, además de licencias permisivas para compartir.

Se declara un modo de trabajo que permite lo popular: los procesos de diseño y fabricación se documentan y publican mientras ocurren, los esquemáticoss, las placas, los códigos y la documentación se escriben en español y se comparten con licencias permisivas, y cada instrumento se piensa para ser reparado, clonado y modificado por otras personas. Dentro de este modo de trabajo, los costos se asumen como parte constitutiva de lo popular: se privilegian componentes asequibles, sin escasez y disponibles localmente, de modo que el precio de un sintetizador no sea una barrera para construirlo, tocarlo o repararlo.

Como marco conceptual para definir e incentivar la popusintesíntesis, se incluye la investigación realizada en el campo de la luthería de instrumentos sonoros de los siglos XX y XXI, con un énfasis en las infraestructuras técnicas y comerciales que los hicieron posibles, además de las metáforas, historia y contexto cultural que las enmarcan.

A nivel de artefactos técnicos que acompañan esta tesis, se incluyen diseño y fabricación de módulos y máquinas que operan como bloques fundamentales de síntesis sonora, para poder remezclar y con ellos construir distintos sintetizadores, y también proponer a partir de ellos derivas y adaptaciones a otros medios como los visuales o textuales, pudiendo así expandirse la popusintesíntesis a distintos medios.

% en borradores iniciales, la extensión sugerida era de 1 página

\end{document}